% ============================================================================
%  chapter6.tex
%  Signals: Asynchronous Events and Their Handlers
% ============================================================================

\chapter{Signals: Asynchronous Events and Their Handlers}

\section{What Is a Signal?}

Everything so far has been synchronous: your code ran, made a system
call, and waited for the answer. Signals are the operating system
breaking that rule. A \emph{signal} is a small, asynchronous
notification delivered to a process. It can arrive at \emph{any}
moment---between any two of your instructions---and it is not requested
by your code. It simply happens to you.

The kernel generates signals for a dizzying variety of reasons:

\begin{verbatim}
   WHO OR WHAT FIRES   SIGNAL        WHY
   ----------------   ----------    ----------------------------------
   you press Ctrl-C    SIGINT        the terminal driver asks you to stop
   you press Ctrl-\    SIGQUIT       stop and leave a core dump
   kill command        SIGTERM       polite shutdown request
   kill -9 command     SIGKILL       unignorable kill
   your program derefs SIGSEGV       invalid memory access (ch.1's abyss)
   a timer expires     SIGALRM       alarm() went off
   a child exits       SIGCHLD       "your child is done; go collect it"
   writing a dead pipe SIGPIPE       the reader of a pipe has gone away
   stop command        SIGSTOP       freeze this process (uncatchable)
   you type a TSTP     SIGTSTP       Ctrl-Z: stop, but catchable
\end{verbatim}

\begin{verbatim}
   Time flows left to right. Signals arrive OUT OF BAND:

     your code:  ----instr--instr--instr--instr--instr--instr--instr---
                              |                  |
                              v                  v
                          SIGINT           SIGSEGV
                          (Ctrl-C)          (bug!)

     The handler runs, returns, and your code resumes at the very next
     instruction. From your code's point of view, a signal is a hidden
     function call injected by the kernel.
\end{verbatim}

A signal is really just a small integer (1--31 on Linux, with a range
of real-time signals beyond that). All the meaning lives in the
\emph{disposition} the process has registered for it and in the default
action the kernel would take if the process has not.

\begin{osconcept}
Only the kernel can deliver signals, because only the kernel can
\emph{interrupt a running process}. A process is a passive guest of the
CPU; it cannot pause itself to receive a notification. The kernel, which
controls the CPU, must notice events (keys, timers, child deaths) and
inject them into the victim process's execution. Asynchronous event
delivery is thus an operating system problem by construction: the OS is
the only actor with the power and the global view to do it.
\end{osconcept}

\section{The Default Actions}

When a signal arrives and the process has no custom handler, the kernel
applies one of a small set of \emph{default actions}. There are exactly
five of them, and every signal's default is one of these:

\begin{verbatim}
   default action     meaning                          examples
   ----------------   -------------------------------  ---------
   terminate          kill the process                SIGHUP SIGINT SIGTERM
   core dump          terminate + save core image     SIGSEGV SIGQUIT SIGABRT
   stop               suspend the process             SIGSTOP SIGTSTP
   continue           resume a stopped process        SIGCONT
   ignore             do nothing                      SIGCHLD SIGURG

   So "Ctrl-C kills my program" is not a rule of the language; it is
   the DEFAULT action for SIGINT, which is why Ctrl-C works the same
   in every program whether or not it knows about signals.
\end{verbatim}

Two signals have no handler, period. \code{SIGKILL} and
\code{SIGSTOP} are the kernel's own kill switch: they cannot be caught,
blocked, or ignored. If any process could block \code{SIGKILL}, a
runaway program could make itself immortal; the OS would lose its
ultimate authority. Their existence is the sign that signals are not a
programmer convenience but an operating system mechanism.

\begin{osconcept}
The default-action table is operating system policy. The kernel must
decide, for every possible event, what happens to a process that has
not expressed an opinion---and that decision must hold even when the
process is broken (a dereferenced \code{NULL}, an illegal instruction).
Core dumps, process reaping, and the uncatchable \code{SIGKILL} are all
designed at the OS level to keep the whole system rescuable, not merely
each individual process.
\end{osconcept}

\section{The Three Dispositions}

For every signal except \code{SIGKILL} and \code{SIGSTOP}, a process
chooses one of three \emph{dispositions}: let the kernel use the
default, ignore the signal outright, or run a function of your own.
That function is the \emph{signal handler}.

\begin{verbatim}
   disposition          what the kernel does when the signal arrives
   ------------------   ----------------------------------------------
   SIG_DFL  (default)   apply the default action from the table above
   SIG_IGN  (ignore)    discard the signal, as if it never happened
   handler  (custom)    call your function, then resume your code

   Default is the starting state of every signal. You override it.
\end{verbatim}

The modern, portable way to install a disposition is
\code{sigaction}, which tells the kernel everything about your handler
in one structure. (The older \code{signal} call behaves differently on
different Unixes; in this book we use \code{sigaction} and forget
\code{signal} exists.)

\begin{lstlisting}[caption={Installing a handler with sigaction.}]
#include <stdio.h>
#include <signal.h>
#include <unistd.h>

/* Runs when SIGINT arrives. "sig" tells you which signal fired. */
static void handle_sigint(int sig) {
    write(1, "interrupted!\n", 13);
}

int main(void) {
    struct sigaction sa;

    sa.sa_handler = handle_sigint;
    sigemptyset(&sa.sa_mask);     /* nothing extra blocked during it */
    sa.sa_flags = 0;

    if (sigaction(SIGINT, &sa, NULL) < 0) {
        perror("sigaction");
        return 1;
    }

    printf("press Ctrl-C (or send SIGINT) ...\n");
    while (1) {
        sleep(1);
    }
    return 0;
}
\end{lstlisting}

\begin{verbatim}
   What just happened:

     struct sigaction sa;      describes the desired behavior
     sa.sa_handler            the function to run
     sa.sa_mask               signals to block WHILE the handler runs
     sa.sa_flags              options (more below)

     sigaction(SIGINT, &sa, NULL)
         ^signal   ^new behavior   ^old behavior (NULL = don't care)

   When Ctrl-C arrives, the kernel saves your state, runs
   handle_sigint(2), restores your state, and your loop continues
   exactly where it was interrupted.
\end{verbatim}

Notice the handler writes with \code{write(2)}, not \code{printf}.
That is not a stylistic tic; it is a hard rule, and we will justify it
in Section~6.6.

\begin{osconcept}
Dispositions are kernel state attached to the process, not program
state. The kernel must remember, for each of dozens of signals, what to
do on delivery, and it must honor that even when the process is deep in
some library you do not control. Registering a handler, validating it,
and switching behavior at the exact moment of delivery is kernel work;
your \code{sigaction} call is just writing an entry in a table the
kernel maintains and consults.
\end{osconcept}

\section{Blocking, Pending, and Delivered}

A signal does not necessarily get delivered the instant it is
generated. The kernel runs a small state machine for each signal of
each process, with three states:

\begin{verbatim}
   GENERATED          BLOCKED?            DELIVERED
   (event happens)    process asked       kernel runs the handler
                      to delay it?        (or default action)
        |                   |
        +----> pending -----+----> if unblocked, deliver now
        |     (waiting)            if blocked, keep pending
        |
        +-----------------+-----------> delivered immediately

   - A BLOCKED signal is not lost: it sits "pending" and is delivered
     the moment the process unblocks it.
   - Standard signals COALESCE: ten SIGINTs generated while blocked
     are delivered as ONE. The kernel remembers "a SIGINT is pending",
     not "ten SIGINTs are pending".
   - A signal that arrives while its own handler is running may be
     delivered again (or blocked, per sa.sa_mask).
\end{verbatim}

Blocking is done with a \emph{signal set}, manipulated through
\code{sigset\_t} and friends:

\begin{lstlisting}[caption={Block, then unblock, a signal.}]
#include <stdio.h>
#include <stdlib.h>
#include <signal.h>
#include <unistd.h>

static void handle_sigint(int sig) {
    write(1, "SIGINT handled\n", 15);
}

int main(void) {
    struct sigaction sa;
    sigset_t block_set;
    sigset_t old_set;
    int i;

    sa.sa_handler = handle_sigint;
    sigemptyset(&sa.sa_mask);
    sa.sa_flags = 0;
    sigaction(SIGINT, &sa, NULL);

    /* 1. Block SIGINT: it will now go pending instead of being delivered. */
    sigemptyset(&block_set);
    sigaddset(&block_set, SIGINT);
    sigprocmask(SIG_BLOCK, &block_set, &old_set);

    printf("SIGINT is now blocked. Send one (Ctrl-C in another terminal,\n");
    printf("or `kill -INT %d` from a shell). Sleeping 3 seconds...\n", getpid());
    sleep(3);

    /* 2. Unblock: any SIGINT that arrived is delivered NOW. */
    printf("unblocking...\n");
    sigprocmask(SIG_SETMASK, &old_set, NULL);

    sleep(2);
    return 0;
}
\end{lstlisting}

\begin{verbatim}
   The blocked interval:

     you send SIGINT           3 seconds later you unblock
        |                             |
        v                             v
   +---------+  blocked,       +-----+  delivered now, handler runs
   | pending |  then one        |     |
   | SIGINT  |  unblock --->    |  !  |----->  write("SIGINT handled")
   +---------+                  +-----+
\end{verbatim}

\begin{ruleofthumb}[Block the signal during its own handler]
Unless you want a handler to be re-entered, add the signal you are
handling to \code{sa.sa\_mask}. Simpler still: make handlers set a
\emph{flag} and do all real work in the main loop (Section~6.7); then
re-entrancy stops being a question you have to answer.
\end{ruleofthumb}

\begin{osconcept}
The pending/blocked/delivered state machine lives entirely in the
kernel. Only the kernel can decide when a process is at a safe point to
run a handler, and only the kernel can hold a signal in limbo without
losing it. A process cannot "pause" its own execution safely---it has
no authority over the scheduler. The signal lifecycle is therefore
kernel bookkeeping, exposed to you as three tiny calls.
\end{osconcept}

\section{What You May Do in a Handler}

Here is the discipline that keeps signal programming honest: a handler
runs \emph{asynchronously}, between two arbitrary instructions of your
main program. If the handler calls a function that your main program is
midway through using---\code{malloc}, for example---the two can corrupt
each other. A handler may therefore call only \emph{async-signal-safe}
functions, a short list blessed by POSIX. \code{printf}, \code{malloc},
\code{strcpy}, and \code{free} are \emph{not} on it. \code{write},
\code{read}, \code{waitpid}, \code{kill}, \code{\_exit}, and
\code{sigaction} are.

\begin{verbatim}
   ASYNC-SIGNAL-SAFE (safe in a handler)     NOT safe in a handler
   -------------------------------          ----------------------
   write(2), read(2)                         printf, sprintf, puts
   _exit(2)                                  exit, malloc, free
   kill(2), raise(3), waitpid(2)             strcpy, strlen, memcpy
   sigaction, sigprocmask                    fopen, fclose, fflush
   alarm(2), getpid(2)                       any code that locks
\end{verbatim}

The consequence is that a handler usually does almost nothing. The
canonical pattern is a \emph{flag}: the handler sets a variable, the
main loop checks it. The variable must be \code{volatile} (so the
compiler does not cache it in a register across the async change) and
of type \code{sig\_atomic\_t} (an integer type guaranteed to be
readable and writable in one atomic step).

\begin{lstlisting}[caption={The flag pattern: handler sets, main loop acts.}]
#include <stdio.h>
#include <signal.h>
#include <unistd.h>

/* Set by the handler, read by main. volatile sig_atomic_t is the
   one type guaranteed safe for this two-way traffic.              */
static volatile sig_atomic_t got_sigint = 0;

static void handle_sigint(int sig) {
    got_sigint = 1;       /* the ONLY thing the handler does */
}

int main(void) {
    struct sigaction sa;
    int ticks;

    sa.sa_handler = handle_sigint;
    sigemptyset(&sa.sa_mask);
    sa.sa_flags = 0;
    sigaction(SIGINT, &sa, NULL);

    ticks = 0;
    while (got_sigint == 0) {
        ticks = ticks + 1;
        printf("tick %d\n", ticks);
        sleep(1);
    }

    printf("SIGINT received; the loop noticed within a second.\n");
    return 0;
}
\end{lstlisting}

\begin{verbatim}
   The division of labor:

     signal arrives ----> handler runs:   got_sigint = 1;
     (kernel thread of execution)         (safe: one atomic store)

     main loop       ----> checks:        while (got_sigint == 0)
     (your thread of execution)           ... does the real work

   The handler and the main loop are like two coroutines sharing one
   int. Everything that is complicated--printing, allocating, closing
   files--belongs in the main loop, where the whole library is safe.
\end{verbatim}

\begin{osconcept}
The async-signal-safe list exists because the kernel delivers signals
at moments it alone chooses, without coordinating with your libraries.
POSIX (the OS standard) must therefore guarantee that a specific subset
of functions never internally allocates, never locks, and never needs
your program's state to be consistent. Whether a function may be called
from a handler is an operating-system-contract decision, not a language
decision. It is the kernel's interrupt model that forces it.
\end{osconcept}

\section{SIGCHLD: Reaping Children Asynchronously}

Chapter~5 ended with a parent that calls \code{waitpid} and blocks
until its child finishes. But what if the parent is busy serving
requests and children finish at unpredictable moments? It cannot sit in
\code{waitpid}. The answer is \code{SIGCHLD}: the kernel sends it to a
parent \emph{whenever} a child changes state---exits, stops, or
continues. The parent installs a \code{SIGCHLD} handler, and the handler
reaps whatever children have died since the last sweep.

The handler must loop: several children may have exited, and
\code{SIGCHLD} \emph{coalesces}, so one signal may stand for many dead
children. The loop drains the pipe:

\begin{lstlisting}[caption={Reaping every dead child with WNOHANG.}]
#include <stdio.h>
#include <stdlib.h>
#include <signal.h>
#include <sys/types.h>
#include <sys/wait.h>
#include <unistd.h>

static volatile sig_atomic_t children_changed = 0;

static void handle_child(int sig) {
    children_changed = 1;
}

/* Called from the MAIN loop, never from the handler itself. */
void reap_children(void) {
    pid_t done;
    int status;

    while (1) {
        done = waitpid(-1, &status, WNOHANG);
        if (done == 0) {
            break;          /* no more children have exited */
        }
        if (done < 0) {
            break;          /* ECHILD: no children left at all */
        }
        if (WIFEXITED(status)) {
            printf("child %d exited with status %d\n",
                   (int) done, WEXITSTATUS(status));
        }
    }
}

int main(void) {
    struct sigaction sa;
    int i;

    sa.sa_handler = handle_child;
    sigemptyset(&sa.sa_mask);
    sa.sa_flags = 0;
    sigaction(SIGCHLD, &sa, NULL);

    /* Spawn three short-lived children. */
    for (i = 0; i < 3; i++) {
        pid_t pid;

        pid = fork();
        if (pid == 0) {
            exit(10 + i);
        }
    }

    /* Busy loop: do work, and check whether reaping is due. */
    for (i = 0; i < 10; i++) {
        sleep(1);
        if (children_changed) {
            children_changed = 0;
            reap_children();
        }
    }

    return 0;
}
\end{lstlisting}

\begin{verbatim}
   What SIGCHLD reaping prevents:

      children exit            zombie count
      -------------            ------------
      3 children finish        the handler's waitpid(-1, WNOHANG)
      (3 SIGCHLD coalesce      loop reaps ALL THREE in one call,
       into possibly 1 signal) so no zombie survives more than a
                               signal delivery.

   WITHOUT the loop (just one waitpid in the handler): two children
   stay zombies until the next SIGCHLD or the parent exits. WITH it,
   the drain loop guarantees every corpse is collected.
\end{verbatim}

\begin{ruleofthumb}[WNOHANG is non-negotiable here]
Inside a \code{SIGCHLD} handler, you must use \code{waitpid(-1,
\&status, WNOHANG)} and loop. A plain \code{waitpid(pid,...)} would
block if no child has exited, freezing the process inside the handler
forever. \code{WNOHANG} means ``don't wait; tell me right now what's
ready,'' which is exactly the contract an asynchronous reaper needs.
\end{ruleofthumb}

\begin{osconcept}
Zombie reaping, introduced in Chapter~5, becomes an event-driven
problem here, and the OS solves it by \emph{fusing} the two mechanisms:
child death generates \code{SIGCHLD} automatically, and \code{waitpid}
collects the result. The kernel chooses to deliver one signal per state
change (coalescing repeats), and it is what keeps the process tree
clean without the parent babysitting every child. Connecting process
lifecycle events to the signal mechanism is a deliberate piece of
operating system design.
\end{osconcept}

\section{Sending Signals: kill and raise}

Signals flow in both directions. \code{kill(pid, sig)} sends a signal
to another process (given permission); \code{raise(sig)} sends one to
yourself; \code{kill(0, sig)} sends to your whole process group. The
permission rule is the kernel's: you may signal a process you own, or
one owned by the same user (or root can signal anyone).

\begin{lstlisting}[caption={Sending signals, three ways.}]
#include <stdio.h>
#include <stdlib.h>
#include <signal.h>
#include <unistd.h>

int main(void) {
    pid_t pid;

    pid = fork();
    if (pid < 0) {
        perror("fork");
        exit(1);
    }

    if (pid == 0) {
        /* Child: announce itself, then idle forever. */
        printf("child %d sleeping...\n", getpid());
        while (1) {
            sleep(1);
        }
    }

    /* Parent: give the child a moment, then say goodnight. */
    sleep(1);
    printf("sending SIGTERM to child %d\n", (int) pid);
    if (kill(pid, SIGTERM) < 0) {
        perror("kill");
        exit(1);
    }

    sleep(1);
    return 0;
}
\end{lstlisting}

\begin{verbatim}
   kill(pid, 0) is a useful trick: it sends NO signal, but succeeds if
   the process exists and you may signal it. "kill -0" is a test of
   existence and permission, widely used in scripts.

   raise(SIGTERM) == kill(getpid(), SIGTERM)   -- signal yourself.
   alarm(5)      : the kernel arranges to send SIGALRM in 5 seconds.
   pause()       : block until ANY signal arrives.
\end{verbatim}

\begin{osconcept}
Permission checking on \code{kill} is kernel work: before delivering,
the kernel verifies the sender owns the target (or is root), because
letting arbitrary processes terminate each other would be a denial of
service. The \emph{capability} to signal someone is a form of
authority the kernel hands out and enforces. \code{kill(pid, 0)}
exists only because that same authority needs a non-destructive test.
\end{osconcept}

\section{alarm: A Signal in the Future}

\code{alarm(n)} asks the kernel to send \code{SIGALRM} to your process
in \code{n} seconds---exactly once. It is the classic way to enforce
timeouts on operations that might hang, because \code{SIGALRM}
interrupts the current system call. Its default action is termination,
so installing a handler (or using it with \code{sigaction}) turns a
hang into a manageable event.

\begin{lstlisting}[caption={A timeout guard around an I/O read.}]
#include <stdio.h>
#include <stdlib.h>
#include <signal.h>
#include <unistd.h>
#include <errno.h>

static volatile sig_atomic_t timed_out = 0;

static void handle_alarm(int sig) {
    timed_out = 1;
}

int main(void) {
    struct sigaction sa;
    char buf[64];
    ssize_t got;

    sa.sa_handler = handle_alarm;
    sigemptyset(&sa.sa_mask);
    sa.sa_flags = 0;
    sigaction(SIGALRM, &sa, NULL);

    printf("you have 5 seconds to type a line\n");

    alarm(5);
    got = read(0, buf, sizeof(buf));     /* may be interrupted by SIGALRM */
    alarm(0);                            /* cancel the timer */

    if (timed_out) {
        printf("\ntimeout; giving up\n");
        return 1;
    }
    if (got < 0) {
        perror("read");
        return 1;
    }

    printf("you typed %zd bytes\n", got);
    return 0;
}
\end{lstlisting}

\begin{verbatim}
   The timeline when the user types nothing:

     t=0  alarm(5)  begins
     t=5  SIGALRM delivered  -->  handle_alarm() sets timed_out
     t=5  read() returns -1, errno = EINTR   (see next section)
     main sees timed_out == 1 and gives up.

   alarm(0) cancels a pending alarm. Real applications time out
   network reads, locks, and whole requests this way--the signal is
   the kernel's ticking clock.
\end{verbatim}

\section{EINTR: When a Signal Interrupts a System Call}

Here is the subtle interaction that causes the classic bug in signal
programming. While your process is blocked inside a \emph{slow}
system call (\code{read} on a pipe or terminal, \code{waitpid},
\code{sleep}), a signal arrives. The kernel must run the handler. The
only safe way is to make the system call return \emph{prematurely}, so
the handler can run in your context. The call returns \code{-1} with
\code{errno = EINTR}, ``interrupted by signal.''

\begin{verbatim}
   blocked in read()               SIGTERM arrives
        |                                |
        |  read returns -1,             v
        |  errno = EINTR          handler runs
        |                                |
        +----- your code: is EINTR an   |
              error or a "retry"?  <-----+
              It is a RETRY: the operation did not fail, it was
              postponed. Standard practice:

                 while (got < 0 && errno == EINTR) {
                     got = read(fd, buf, n);
                 }
\end{verbatim}

Two ways to cope. The manual way wraps every slow call in an
\code{EINTR} retry loop. The automatic way sets \code{SA\_RESTART} in
\code{sa.sa\_flags}, asking the kernel to restart interrupted calls
itself---but note that \code{read} on a pipe, \code{waitpid}, and a few
others are \emph{not} restartable by it, so the manual loop is the
portable truth.

\begin{lstlisting}[caption={The universal EINTR retry loop.}]
#include <unistd.h>
#include <errno.h>

/* Read up to n bytes, restarting if a signal interrupts us. */
ssize_t read_eintr(int fd, void *buf, size_t n) {
    ssize_t got;

    while (1) {
        got = read(fd, buf, n);
        if (got < 0 && errno == EINTR) {
            continue;      /* a signal arrived; retry, don't give up */
        }
        return got;        /* real result: n, k<n, 0 (EOF), or an error */
    }
}
\end{lstlisting}

\begin{osconcept}
The kernel's decision to interrupt slow system calls is what makes
asynchronous delivery \emph{possible at all}: a handler cannot run
inside a blocked call, so the kernel must unwind the call, deliver, and
let your code decide what to do next. \code{EINTR} is therefore not a
bug in your code; it is the visible trace of the kernel's interrupt
mechanism. Designing which calls restart automatically and which do
not is an operating system judgment call, and \code{SA\_RESTART} is the
kernel exposing that choice to you.
\end{osconcept}

\section{Practical: A Graceful Shutdown Server}

Let us assemble the chapter into the pattern that real daemons use:
handle \code{SIGTERM}/\code{SIGINT} by setting a flag, check it in the
main loop, and reap children via \code{SIGCHLD}. The handlers do almost
nothing; all work happens in the main loop where the full standard
library is safe.

\begin{lstlisting}[caption={server.c: flag-based shutdown plus asynchronous reaping.}]
#include <stdio.h>
#include <stdlib.h>
#include <signal.h>
#include <unistd.h>
#include <sys/types.h>
#include <sys/wait.h>

static volatile sig_atomic_t g_running = 1;
static volatile sig_atomic_t g_children = 0;

static void handle_shutdown(int sig) {
    g_running = 0;          /* ask the main loop to finish */
}

static void handle_child(int sig) {
    g_children = 1;         /* ask the main loop to reap   */
}

void reap_children(void) {
    pid_t done;
    int status;

    while (1) {
        done = waitpid(-1, &status, WNOHANG);
        if (done <= 0) {
            break;
        }
        if (WIFEXITED(status)) {
            printf("reaped child %d (status %d)\n",
                   (int) done, WEXITSTATUS(status));
        } else if (WIFSIGNALED(status)) {
            printf("reaped child %d (signal %d)\n",
                   (int) done, WTERMSIG(status));
        }
    }
}

int main(void) {
    struct sigaction sa;

    /* One handler for both termination signals. */
    sa.sa_handler = handle_shutdown;
    sigemptyset(&sa.sa_mask);
    sa.sa_flags = 0;
    sigaction(SIGINT, &sa, NULL);
    sigaction(SIGTERM, &sa, NULL);

    sa.sa_handler = handle_child;
    sigemptyset(&sa.sa_mask);
    sa.sa_flags = 0;
    sigaction(SIGCHLD, &sa, NULL);

    printf("server running (pid %d). Try: kill -TERM %d\n",
           getpid(), getpid());

    while (g_running) {
        if (g_children) {
            g_children = 0;
            reap_children();
        }
        /* In a real server: accept a connection, fork a worker... */
        sleep(1);
    }

    printf("shutdown complete; resources cleaned up\n");
    return 0;
}
\end{lstlisting}

\begin{verbatim}
   What happens end to end:

     kill -TERM <pid>  --->  SIGTERM  --->  handle_shutdown
                                              |  g_running = 0
     the main loop's next check of g_running  |
     sees 0, exits the loop, cleans up,       |
     returns from main.  <-------------------+

     A child exits  ----->  SIGCHLD  --->  handle_child
                                              |  g_children = 1
     the main loop reaps every zombie with    |
     waitpid(-1, WNOHANG).  <----------------+

   Both handlers are two assignments. The main loop is the only place
   that prints, allocates, or calls anything complex. That is the whole
   art of signal programming.
\end{verbatim}

\begin{researchnote}[Why the flag-and-loop pattern wins]
The graceful-shutdown server you just built is \emph{event-driven}: one
loop, a few flags, and all slow work done inline. This is not merely a
style preference; it is a design conclusion with a famous champion.
John Ousterhout, the creator of Tcl/Tk and a leading systems
researcher, argued in a widely-cited essay that a single event-driven
loop is easier to make correct than a design where many concurrent
activities each block on their own.

\begin{verbatim}
   The essay's central claim (paraphrased):

       "Concurrency via threads is hard to program correctly, and for
        most applications an event-driven design -- one loop that
        checks state and handles events as they arrive -- is simpler
        and more robust."

   Translate to this chapter: our main loop checks g_running and
   g_children instead of blocking; the signal handlers only set flags.
   That is event-driven thinking applied to asynchronous events.
\end{verbatim}

Why this matters beyond signals: it is the same architecture we will
use for network servers in Chapter~17, where one loop multiplexes
thousands of connections. The claim is contested (some argue threads
scale better), but the pattern is proven and it is what our code uses.

\paperref{J.~Ousterhout, ``Why Threads Are A Bad Idea (for most
purposes),'' invited talk, USENIX 1996}{https://web.stanford.edu/~ouster/cgi-bin/papers/threads.pdf}
\end{researchnote}

\section{Structures in Bytes}

A signal is a small integer (1--31 standard, plus 34--64 real-time on
Linux), so the kernel's per-process signal state is mostly bits. The
\code{sigset\_t} you pass to \code{sigprocmask} is exactly a bit array:

\begin{verbatim}
   sigset_t on Linux/x86-64: 128 bytes = 1024 bits, one bit per signal.

   byte:   0        1        2        ...   127
          +--------+--------+--------+-- --+--------+
          |        |        |        |     |        |   bit i = signal i+1
          +--------+--------+--------+-- --+--------+
          SIGINT (2)  -> bit 1 of byte 0
          SIGKILL (9) -> bit 8: byte 1, bit 0
          SIGTERM(15) -> bit 14: byte 1, bit 6

   struct sigaction on glibc/x86-64: 152 bytes
   offset  size  field
   0       8     sa_handler      (or sa_sigaction: same storage)
   8       128   sa_mask         (a sigset_t)
   136     4     sa_flags        (SA_RESTART, SA_SIGINFO, ...)
   144     8     sa_restorer     (legacy; you never set it)
\end{verbatim}

(The exact sizes are an ABI detail; the shape is what matters: one
pointer, one bit-array, one int, one pointer.) You never touch these
bytes directly---\code{sigemptyset} and \code{sigaddset} manage the
bit array, \code{sigaction} fills the struct---but it is useful to
know that blocking a signal is just setting a bit.

\section{Summary}

\begin{verbatim}
   The signal model in one picture:

      event ------> GENERATED
                       |      blocked? yes ----> PENDING (coalesced)
                       |      blocked? no       | (delivered on unblock)
                       v                        v
                  DELIVERED ----> default action or your handler
                       |
                       +--- handler may set flags (volatile sig_atomic_t)
                       +--- handler may call only async-signal-safe fns
                       +--- slow syscalls return EINTR; retry or SA_RESTART

   The rules that keep signal code correct:
     1. install with sigaction, not signal;
     2. handlers set flags; the main loop does the work;
     3. flags are volatile sig_atomic_t;
     4. never call malloc/printf in a handler; use write/_exit;
     5. reap SIGCHLD with a WNOHANG loop;
     6. wrap slow syscalls against EINTR;
     7. SIGKILL and SIGSTOP are not yours to catch.
\end{verbatim}

\begin{osconcept}
Signals are the oldest asynchronous mechanism in Unix, and they remain
the OS's primary way of injecting events into running processes. The
whole chapter is a tour of operating system design: privileged event
delivery, per-process dispositions, a kernel-side state machine
(pending/blocked/delivered), an async-safety contract, and the
interaction of signals with blocking system calls. Master it and you
can build daemons, supervisors, and servers that behave correctly under
the chaos of the real machine.
\end{osconcept}

In the next chapter we give signals a channel: inter-process
communication through pipes, which let one process pour its output into
another---the mechanism behind every command pipeline you have ever
typed.

